\section{Detaillierte Analyse ausgewählter Module}
Um die technische Tiefe zu verdeutlichen, wird hier praktisch die Implementierung sicherheitsrelevanter und komplexer Logik dargestellt.

\subsection{Authentifizierung: Secure Cookies und Hashing}
Die Sicherheit der Benutzerdaten hat oberste Priorität. Der \texttt{AuthService} implementiert daher bewährte Industriestandards.

\subsubsection{Passwort-Hashing (BCrypt)}
Passwörter werden niemals im Klartext gespeichert. Stattdessen wird der \texttt{BCrypt}-Algorithmus verwendet, der durch seinen konfigurierbaren "Work Factor" (Kostenfaktor) resistent gegen Brute-Force- und Rainbow-Table-Angriffe ist.
\begin{lstlisting}[language={[Sharp]C}, caption={Verifizierung im AuthService}]
// user.PasswordHash kommt aus der DB, passwordDto.Password vom User
bool isPasswordValid = BCrypt.Net.BCrypt.Verify(passwordDto.Password, user.PasswordHash);
if (!isPasswordValid)
{
    throw new UnauthorizedAccessException("Invalid credentials");
}
\end{lstlisting}

\subsubsection{Cookie-Konstruktion}
Nach erfolgreichem Login wird kein JWT im LocalStorage (Frontend) abgelegt (XSS-Gefahr), sondern ein \textbf{HttpOnly Cookie} gesetzt.
\begin{itemize}
    \item \textbf{HttpOnly:} Das Cookie kann nicht per JavaScript ausgelesen werden.
    \item \textbf{Secure:} Wird nur über HTTPS übertragen.
    \item \textbf{SameSite=Strict:} Das Cookie wird nur an den eigenen Origin gesendet.
\end{itemize}

\subsection{Ingestion-Pipeline: HNSW Indexierung}
Wie in Abschnitt 6 beschrieben, nutzt das System einen HNSW-Index. Die technische Herausforderung liegt hier in der Wahl der passenden Index-Strategie für die Vektoren innerhalb von \texttt{pgvector}. Code-seitig wird dies in der \texttt{OnModelCreating}-Methode des \texttt{ApiDbContext} gelöst, um sicherzustellen, dass PostgreSQL den Index korrekt für L2-Distanz-Berechnungen optimiert. Dies ist ein direktes Mapping von C\#-Code auf datenbankspezifische Features.
